% Lean module: MonadFormal.Witnesses
\section{Separating counterexamples}\label{sec:separating-witnesses}

Several distinctions that look harmless in isolation fail when
strengthened. The examples below separate pairwise from joint
agreement, macro labels from implemented interventions, approximate
closeness from equivalence, and reliable recording from thresholded
detection.

\begin{theorem}[Pairwise agreement does not determine the joint law]
\label{thm:parity-marginals}
Let $P$ be uniform on $\{000,011,101,110\}$ and $Q$ uniform on
$\{001,010,100,111\}$. Every two-coordinate marginal agrees between
$P$ and $Q$, but parity is always zero under $P$ and always one under $Q$.
\end{theorem}
\begin{proof}
Each support contains four states. For every two distinct coordinates,
each of $00,01,10,11$ appears exactly once in each support, giving
probability $1/4$. Direct evaluation of exclusive-or gives zero on
every member of the first support and one on every member of the second.
\leanproof{MonadFormal.Witnesses.equal_support_sizes,MonadFormal.Witnesses.every_pair_marginal_equal,MonadFormal.Witnesses.opposite_joint_parity}
\end{proof}
Pairwise validation is therefore not joint validation. Resource admission
is a separate issue even before statistical inference: three measurements
of cost four have jointly inadmissible cost under budget ten although any
pair is affordable.
\begin{proof}
The costs are $4+4=8\le10<12=4+4+4$.
\leanproof{MonadFormal.Witnesses.pair_budget_not_joint_budget}
\end{proof}

\begin{theorem}[A macro intervention needs an implementation]
\label{thm:macro-implementation}
For microstates $(x_1,x_2)\in\mathbb N^2$, let
$a(x_1,x_2)=x_1+x_2$ and $b(x_1,x_2)=x_1+2x_2$.
Setting the macro value $a=1$ via $(1,0)$ yields $b=1$, whereas
setting it via $(0,1)$ yields $b=2$. Consequently no deterministic
response function $h$ satisfies $b=h\circ a$ on all microstates.
\end{theorem}
\begin{proof}
Both microstates have $a=1$ and their $b$ values differ. A function
of the macro value alone would assign the same $b$ value to both.
\leanproof{MonadFormal.Witnesses.same_macro_input_different_output,MonadFormal.Witnesses.no_macro_response}
\end{proof}
This does not forbid a macro model with a specified implementation
distribution or additional mechanism constraints. It shows why the macro
label alone is insufficient to determine the intervention's response.

\begin{lemma}[Binary total variation]\label{lem:binary-tv}
For Bernoulli parameters $p,q\in[0,1]$, total variation is $|p-q|$.
\end{lemma}
\begin{proof}
For a test $f:\{0,1\}\to[0,1]$, the expectation difference is
$(p-q)(f(1)-f(0))$, whose absolute value is at most $|p-q|$.
The indicator of one attains this bound. Equivalently, the two-atom
expression $\frac12(|p-q|+|(1-p)-(1-q)|)$ has equal absolute-value terms.
\leanproof{MonadFormal.Witnesses.bernoulli,MonadFormal.Witnesses.bernoulli_tv,MonadFormal.Witnesses.binary_total_variation}
\end{proof}

\begin{example}[Approximate closeness is not equivalence]
\label{ex:approximate-nontransitivity}
At total-variation tolerance $1/2$, Bernoulli parameters zero and $2/5$
are close, as are $2/5$ and $4/5$, but zero and $4/5$ are not.
\end{example}
\begin{proof}
The adjacent distances are $2/5$ and the endpoint distance is $4/5$.
\leanproof{MonadFormal.Witnesses.approximate_closeness_not_transitive,MonadFormal.Witnesses.approximate_tv_not_transitive}
\end{proof}
Thus an approximate MONAD relation cannot automatically be quotiented
as an equivalence relation.

\begin{example}[Reliable but below the detection threshold]
\label{ex:record-threshold}
Suppose the two event-conditioned binary records have probabilities
$4/5$ and $1/5$ of recording one. The identity decision rule has both
error probabilities $1/5<1/2$. Nevertheless their total variation is
$3/5$, below a declared detection threshold $9/10$.
\end{example}
\begin{proof}
The power is $4/5=1-1/5$, the size is $1/5$, and
$|4/5-1/5|=3/5<9/10$.
\leanproof{MonadFormal.Witnesses.recordable_but_undetected}
\end{proof}
